\chapter{69}\label{section-69}

«Ja, ich war bei Giorgias Eltern, ich wollte wissen, wie es ihr geht,»
Valentina sprach gefasst, fast so, als hätte sie Louis’ Anruf erwartet
und sich vorbereitet, «die haben nun das künstliche Koma verlängert. Am
Samstag wurde sie am Rücken operiert. Es sei soweit gut gelaufen, sagen
die Ärzte. In ein paar Tagen soll Giorgia langsam aufgeweckt werden. Die
Salas sind, wie soll ich sagen,» sie suchte nach Worten, «immer noch
fassungslos. Und es ist ungewiss, wie es Giorgia geht, wenn sie
aufwacht. Ihre Mutter hat geweint und die Polizei war einige Male bei
ihnen. Die, die warten darauf, endlich mit Giorgia sprechen zu können,
weil….»

«Verrückt, sie ist also schon bald zwei Wochen nicht ansprechbar?»

Louis wusste nicht, ob ihn die Nachricht beruhigen oder besorgen sollte.

«Ja, darum ist es so kompliziert, zu verstehen, was auf dem
\emph{Roccia} geschehen ist. Die Polizei hat Basti laufen lassen,
mangels Beweisen, sagen sie, aber Giorgias Vater meint…»

Louis unterbrach Valentina erneut.

«Aber was soll Bastiano denn getan haben?»

«Ich weiss kaum mehr, als ich dir bereits gesagt habe. Die Salas halten
sich bedeckt. Es fällt ihnen schwer, Bastiano zu verdächtigen; der war
viele Jahre wie ein Kind für sie, bevor er verschwand, damals. Sie
meinten nur, Bastis Vater sei zu einflussreich und die Polizei kusche,
ach, ich weiss nicht.»

Louis hörte Valentina Luft auspusten.

«Auf jeden Fall scheint der leitende Commissario ziemlich sicher zu
sein, dass es kein Selbstunfall war. Und er soll sich bei den Salas nach
dem Drogenkonsum ihrer Tochter erkundigt haben.»

«Ach.»

Drogen? Die Bilder von Giorgia, die Worte, die sie ihm an den Kopf
geworfen hatte, alles war wieder da, als spiele sich die Szene gerade
eben ab. Waren Drogen die Erklärung für Giorgias Verhalten? Sie hatte
Drogen vehement abgelehnt, selbst Joints waren ihr zuwider. Von Bier und
Wein wurde sie lustig. Betrunken hatte er sie selten erlebt. Giorgia und
Drogen. Das passte nicht zusammen. Louis schluckte.

«Aber was erzähle ich dir, Louis, du weisst doch, wie’s war, oder?»

Louis blieb dabei, er habe Giorgia abstürzen sehen. Und nichts weiter.
Valentina liess ihn mit ihrem Schweigen ahnen, dass er sie enttäuschte.
Für einen Moment war Louis drauf und dran, sie einzuweihen.

«Freitagabend rief mich dieser Sömnez von der Polizei an und lud mich
für Samstag auf die Questur vor, zum zweiten Mal. Sie hätten noch ein
paar Fragen, sagte er.»

«Was wollten sie denn noch von dir hören?»

«Hm, du willst es nicht wissen, Louis.»

«Was denn?»

Louis’ Angst dehnte sich aus.

«Die Salas erzählten offenbar von dir, deiner Beziehung mit Giorgia, dem
Ende. Der Commissario, Mantovani heisst der, schien jedenfalls gut
informiert. Der wusste auch, dass wir befreundet sind, du und ich.»

Louis’ Herz galoppierte.

«Was hast du ihm erzählt?»

«Das Gleiche wie das erste Mal, nur, der hat mich richtig in die Mangel
genommen.»

Valentinas Stimme würde nächstens versagen. Louis hoffte, dass sie nicht
wieder zu weinen anfing.

«Er fragte, warum ich so nervös sei, warum ich seinem Blick ausweiche,
und er vermute, ich wisse mehr. Louis, die Befragung war grauenhaft. Der
wollte von mir wissen, warum du nicht erreichbar bist. Und ich sagte,
dass ich das nicht weiss, dass ich auch gerade keinen Kontakt zu dir
habe, aber, Louis, der hat mir nicht geglaubt, dann liess er mich gehen.
Ich hab ihn angelogen. Das war Scheisse.»

Ihr letzter Satz klang nach. Damit hatte Louis nicht gerechnet. Er war
sprachlos und brauchte erstmal Zeit.

Ihm blieb einzig, sich bei ihr zu bedanken. Mit dem Vorwand, er müsse
dringend weg, vertröstete er sie und verabschiedete sich. Kurz und
knapp. Und mit dem miesen Gefühl, ein schlechter Freund zu sein. Er
hatte Valentina mithineingezogen. Er verlangte von ihr, dass sie ihn
deckte. Und er wusste, dass es falsch war.
